%=============================================================================!
% 16.1  HOW ATOMS ARRANGE THEMSELVES                                          !
%=============================================================================!
\section{How atoms arrange themselves}
\label{sec:ob-rdf}

An orchard and a wood can hold the same number of trees on the same area
and still look different. Standing at one tree of an orchard planted on a
square grid, a walker finds the nearest trees at one spacing, the next at
$\sqrt2$ spacings, then at 2, each at a fixed distance and in a fixed
number. In a wood no two trees grow closer than their trunks allow, the
nearest neighbours lie at roughly one distance without fixed directions,
and farther out the trees thin into an even scatter. Counting the trees in
rings of equal width around one tree, and comparing each count with what
an even scatter would give, tells the two apart without a map. Atoms in a
crystal and in a liquid differ in the same way, and the count in rings is
the radial distribution function.

\subsection*{The radial distribution function}

Around an atom, the number of others whose centres lie between the
distances $r$ and $r + \dd r$ would be the number density $\rho = N/V$
times the volume of that thin shell, $4\pi r^2\,\dd r$, if the atoms were
scattered evenly and independently. The radial distribution function
$g(r)$ is the factor by which the actual mean number differs:
\begin{equation}
   \text{mean number of partners between $r$ and $r + \dd r$}
   = \rho\,g(r)\,4\pi r^2\,\dd r .
   \label{eq:ob-rdf}
\end{equation}
Where $g(r) > 1$ partners crowd; where $g(r) < 1$ they are scarce; and
$g(r) = 0$ inside the distance at which two atoms repel too strongly to
meet. The mean is over every atom as the centre and over every frame of
the run, so $g(r)$ describes the whole liquid or crystal.
Summing \cref{eq:ob-rdf} out to a distance $r$ counts the partners within
$r$ of an atom, the running coordination number
\begin{equation*}
   n_{\mathrm c}(r) = \rho\int_0^r g(s)\,4\pi s^2\,\dd s ,
\end{equation*}
whose value at the end of the first shell of neighbours is the atom's
number of nearest neighbours.

A run gives $g(r)$ as a histogram (\cref{sec:sm-probability}). The
distances are cut into bins of width $\delta r$; in each frame, every pair
closer than the last edge adds one to its bin for each of its two atoms,
since each is the other's partner; and the total over $n$ frames is
divided by the count an even scatter would give, $nN\rho$ times the volume
of the shell between the edges $r_k$ and $r_{k+1}$,
\begin{equation*}
   \tfrac43\pi\bigl(r_{k+1}^3 - r_k^3\bigr)
   = 4\pi r_k^2\,\delta r + 4\pi r_k\,\delta r^2
   + \tfrac43\pi\,\delta r^3 ,
\end{equation*}
by expanding the cube of $r_k + \delta r$. The exact difference of
spheres is used rather than its first term $4\pi r_k^2\,\delta r$, which
misses the shell's volume by a fraction of about $\delta r/r_k$, large near
the origin and the whole of it in the first bin, where $r_k = 0$. In a
periodic cell the separations are reduced to their nearest image
(\cref{sec:md-image}). That finds every pair once only out to half the
cell's smallest width; beyond it a sphere around an atom would hold more
than one copy of some partners while the nearest image counts one, so the
histogram stops there. A partial function $g_{\mathrm{AB}}(r)$, the
partners of kind B around atoms of kind A, takes A atoms as the centres,
B atoms as the partners and $\rho_{\mathrm B} = N_{\mathrm B}/V$ as the
density.

The function \texttt{rdf} of \texttt{mdlab.analysis.structure} does this,
finding the pairs through the cell list of \cref{sec:md-neighbours}:

\begin{code}
def rdf(
    positions: ArrayLike,
    cell: ArrayLike,
    r_max: float,
    n_bins: int,
    centres: ArrayLike | None = None,
    partners: ArrayLike | None = None,
    per_frame: bool = False,
) -> tuple[NDArray, NDArray]:
    frames = _frames(positions)
    h = np.asarray(cell, dtype=float)
    in_a = _chosen(frames.shape[1], centres)
    in_b = _chosen(frames.shape[1], partners)
    edges = np.linspace(0.0, r_max, n_bins + 1)
    ideal = in_a.sum() * in_b.sum() / cell_volume(h) * shell_volumes(edges)
    out = np.empty((len(frames), n_bins))
    for k, r in enumerate(frames):
        i, j, _, dist = cell_list_pairs(r, h, r_max)
        # a pair counts once for each end that is a centre whose partner
        # is the other end
        weight = (in_a[i] & in_b[j]).astype(float) + (in_a[j] & in_b[i])
        out[k] = np.histogram(dist, edges, weights=weight)[0] / ideal
    middle = 0.5 * (edges[1:] + edges[:-1])
    return middle, (out if per_frame else out.mean(axis=0))
\end{code}

\noindent It divides by $N\rho = N^2/V$, as ASE's \texttt{get\_rdf}
does\cite{Larsen2017}, and agrees with it to \num{3e-14} on 50 frames of
liquid argon. One consequence of that choice shows far from any atom.
For atoms placed independently at random, an atom's partners are the
other $N - 1$, spread evenly over the volume $V$, while the divisor
assumes $N$, so $g(r)$ tends to $(N - 1)/N = 1 - 1/N$. For 256 atoms
placed at random in the liquid's cell, 400 times over, $g(r)$ averages
$0.9961 \pm 0.0005$ beyond \SI{2}{\angstrom}, against $1 - 1/256 = 0.9961$
(\cref{fig:ob-rdf}a). A liquid falls short of 1 by less. Each atom clears
a hole around itself that holds fewer partners than an even scatter, and
\cref{sec:ob-sq} shows that the hole takes $1 - S(0)$ of the one partner
missing from $\rho V$, with $S(0)$ about 0.06 here, leaving only $S(0)$
missing from the rest of the cell; $g(r)$ then tends to $1 - S(0)/N$, once
the cell is large enough for the liquid's order to die away within it.
Beyond the sphere of radius $L/2$, in the corners of the 2048-atom cell
of Chapter~15, the partners lie at $0.999967 \pm 0.000005$ of $\rho$,
against $1 - 1/2048 = 0.999512$. In the 256-atom cell, where $g(r)$ still
ripples by a few hundredths at $L/2$, no far value is reached.

\subsection*{A crystal and a liquid}

The fcc crystal of the model argon of Chapters~10 to 15, at its own
lattice constant of \SI{5.26714}{\angstrom} and \SI{21}{\kelvin}, has
$g(r)$ in sharp peaks at the distances of its shells of neighbours,
$(a/\sqrt2)\sqrt m$, for $m = 1$ to 5: \num{3.724}, \num{5.267},
\num{6.451}, \num{7.449} and \SI{8.328}{\angstrom}
(\cref{fig:ob-rdf}b). Between the peaks $g(r)$ is zero, and $n_{\mathrm c}(r)$
climbs in steps of 12, 6, 24 and 12, the numbers of atoms in the first
four shells of the fcc lattice, to 12, 18, 42 and 54 at the gaps between
them (\cref{fig:ob-rdf}c). The peaks have width because the atoms vibrate,
each wandering from its mean position by \SI{0.094}{\angstrom} along each
axis (root mean square), and the first rises to 9.30.

The liquid at \SI{135}{\kelvin} and $\rho = \SI{0.02035}{\per\angstrom
\cubed}$ (Chapter~15's run of \SI{2}{\nano\second}, every
\SI{2}{\pico\second}) has none of this order beyond a few neighbours.
$g(r)$ is zero inside \SI{3}{\angstrom}, rises to a first peak of 2.578 at
\SI{3.684}{\angstrom}, falls to a first minimum of 0.663 at
\SI{5.238}{\angstrom}, and then oscillates about 1 with peaks that weaken
shell by shell. Out to its first minimum the liquid holds 12.00
neighbours on average, as many as the crystal's first shell, though
spread over a wider range of distances, at no fixed directions, and at a
density of \SI{0.02035}{\per\angstrom\cubed} against the crystal's
\SI{0.02737}{\per\angstrom\cubed}.

\begin{figure}[htb]
\centering
\includegraphics{ch16_observables/rdf}
\caption{(a) The radial distribution function of liquid argon at
\SI{135}{\kelvin} (Chapter~15's run of \SI{2}{\nano\second}, every
\SI{2}{\pico\second}) against 256 atoms at random in the same cell
(dashed, from \SI{1.5}{\angstrom}, where its bins hold enough pairs).
(b) The fcc crystal at \SI{21}{\kelvin}, with the distances of its first
five shells (dotted). (c) The running coordination number of the liquid
and the crystal, with the liquid's first minimum (dotted).}
\label{fig:ob-rdf}
\end{figure}

\subsection*{Energy and pressure from the arrangement}

For a pair potential, $g(r)$ holds all that the potential energy and the
virial need. The potential energy is the sum of $\varphi(r_{ij})$ over pairs.
Around one atom, \cref{eq:ob-rdf} puts $\rho g(r)4\pi r^2\dd r$ partners
at the distance $r$, each contributing $\varphi(r)$; summing over the $N$
atoms counts each pair from both ends, so
\begin{equation*}
   \ens{U} = \frac N2\,\rho\int_0^\infty\varphi(r)\,g(r)\,4\pi r^2\,\dd r .
\end{equation*}
The virial of \cref{sec:pr-periodic}, $\mathcal{W} = -\sum_{\text{pairs}}
r_{ij}\varphi'(r_{ij})$, follows by the same count,
\begin{equation*}
   \ens{\mathcal{W}} = -\frac N2\,\rho\int_0^\infty r\varphi'(r)\,g(r)\,
   4\pi r^2\,\dd r ,
\end{equation*}
and with it and the run's mean kinetic energy the pressure $\ens{P} =
(2\ens{K} + \ens{\mathcal{W}})/3V$ of \cref{eq:pr-virial}. On bins of
\SI{0.01}{\angstrom} and the 1000 frames, these give $U/N =
\SI{-50.102}{\milli\electronvolt}$ against the run's average of $(-50.092
\pm 0.018)$\,\si{\milli\electronvolt}, 0.6 errors away, and $P =
\SI{0.0848}{\giga\pascal}$ against $(0.0849 \pm 0.0004)$\,\si{\giga
\pascal}, 0.2 errors away. The bins must be fine: on bins of
\SI{0.05}{\angstrom} the energy misses by 1.9 errors, because $\varphi$
changes steeply across a bin near the core. Agreement here is a check of
the normalisation of $g(r)$, of its histogram and of the run's energy
together.

\subsection*{How long the arrangement takes}

Chapter~15 deferred the convergence of $g(r)$ to this chapter. One frame
already averages over 256 centres, so in bins of \SI{0.05}{\angstrom} its
$g(r)$ at the first peak spreads by only 0.33 from frame to frame, and
frames \SI{500}{\femto\second} apart are nearly independent, their
statistical inefficiency (\cref{sec:cv-correlation}) being 1.04, so
$\tau_{\mathrm{int}} = \SI{259}{\femto\second}$. With the run's length $t_{\mathrm f}$ in \cref{eq:cv-length},
the error of the peak's height is 0.076 after \SI{10}{\pico\second}
(\num{2.9}\% of it), 0.024 after \SI{100}{\pico\second} (0.9\%) and 0.005
after \SI{2}{\nano\second}, and a run of \SI{573}{\pico\second} gives
it to 0.01. Pieces of the run of 5 to \SI{200}{\pico\second} scatter by
0.102 to 0.020 against the predicted 0.107 to 0.017, within the scatter
expected of so few pieces, about a quarter for the 10 pieces of
\SI{200}{\pico\second}, so these error bars can be trusted
(\cref{fig:ob-rdf-error}b). Over the whole curve, the
first \SI{10}{\pico\second} lie within two of their errors of the
\SI{2}{\nano\second} result in 0.95 of the bins (\cref{fig:ob-rdf-error}a).
The first peak of $g(r)$ is fixed to 0.9\% by \SI{100}{\pico\second},
where the pressure and the diffusion coefficient need \SI{1}{\nano\second}
for 0.7\% and 1.0\% (\cref{sec:cv-observables}). Its memory is short:
$\tau_{\mathrm{int}}$ is \SI{0.26}{\pico\second}, against
\SI{0.94}{\pico\second} for the pressure under the same thermostat.

\begin{figure}[htb]
\centering
\includegraphics{ch16_observables/rdf_error}
\caption{(a) The radial distribution function of liquid argon from the
first \SI{10}{\pico\second} of Chapter~15's run, less that from all
\SI{2}{\nano\second}, with twice its standard error from the 20 frames
\SI{500}{\femto\second} apart (band). (b) The standard error of $g$ at its
first peak, in bins of \SI{0.05}{\angstrom}, against the length of run
used: predicted from one frame's
spread and the statistical inefficiency (dashed), and measured as the
scatter of independent pieces (points).}
\label{fig:ob-rdf-error}
\end{figure}

\begin{keyidea}
The radial distribution function $g(r)$ counts the partners at the
distance $r$ from an atom against an even scatter. It is a histogram of
pair distances divided by $N\rho$ times the volume of each shell, out to
half the cell, and tends far away to $1 - 1/N$ for atoms placed at random
and to $1 - S(0)/N$, nearly 1, in a liquid in a large enough cell. Its
peaks give the shells of neighbours, its integral the coordination
number, and with the pair potential the energy and the pressure. Its
memory is short, so it converges several times sooner than the
pressure.
\end{keyidea}

\notebook{16}{grow shells around an atom and watch the histogram become g(r)}

\begin{selfcheck}
A gas of 500 atoms placed independently at random in a fixed cell is
analysed with ASE's
normalisation. What value should $g(r)$ approach far from any atom, and
why does a larger cell bring it closer to 1?
\end{selfcheck}
